\scanpage{142}\mspage{\deleted{2}16}
% archive sheet number: 141
\unsure{indexées|indexés} par ces ensembles, savoir les $X_D^{C_0}$ et les $X_D^{C_1}$ relatifs aux diagrammes $D$ impliquant, soit la flèche $f$, soit la flèche $f'$\ldots.

Or la construction \inserted{canonique} du système inductif ordinal $\unsure{P_\alpha|p_\alpha}$, avec les $i_\alpha:\mathrm{id}_{\uline{\mathrm{Fl}}\,M}\to P_\alpha$, est fonctorielle \unsure{p.v.|p.r.} à la donnée initiale $P_1,i_1$ -- je \deleted{\unsure{les|le|ne}} n'explicite pas le détail de la fonctorialité, bien que je devrais\ldots

3. \uline{Majoration des Cardinaux}

Je suppose à présent que $M$ est muni d'une \uline{filtration cardinale}
\[
(\mathrm{Filt}^{\unsure{\pi|\Pi}})_{\pi\geq\pi_0}\qquad(\pi_0,\pi\text{ cardinaux }\in\unsure{\mathcal U|U})\tag*{(3.1.)}
\]
(SGA 4 I 9.12, p. 149), \deleted{\unsure{avec les}} les $\mathrm{Filt}^{\pi}$ stables par sommes amalgamées simples (chaque fois qu'elles existent), donc aussi par sommes amalgamées infinies indexées par un ens. d'indices de cardinal $\leq\pi$. On suppose de plus que \unsure{tout|tous|les} \unsure{objet|objets} de $\mathrm{Filt}^{\pi_0}$ sont accessibles, et \deleted{soit} on désigne par \textcircled{$\pi_1$} le plus petit cardinal$^*$ tel que $\pi_1\geq\pi_0$ et que les objets de $\mathrm{Filt}^{\pi_0}$ soient $\pi_1$-accessibles. \unsure{Désignant|Désignons}, pour \deleted{$\unsure{\pi}$} \unsure{tt} cardinal, par \deleted{$\unsure{\mathrm{Acc}^{\pi}|\mathrm{Filt}^{\pi}|\mathcal A^{\pi}}$} \inserted{$M_\pi$} le sous-ensemble de $\mathrm{Ob}\,M$ formé des objets $\pi$-accessibles, on aura donc pour \unsure{tt} $\pi\geq\pi_1$ (loc cit. 9.16)
\[
\deleted{\unsure{A}}M_{\pi}^{\deleted{\unsure{\pi|*}}}\subset\mathrm{Filt}^{\pi}\subset M_{\deleted{\unsure{\pi}}\unsure{\pi|\deleted{\pi}}\unsure{\times|\deleted{\times}}(\pi^{\pi_0})}^{\deleted{\unsure{\pi|*}}},\quad(\text{si }\pi\geq\pi_1)\tag*{(3.2)}
\]
